\section{Finding Borders}
\label{sec:cses-finding-borders}

\problemstatement{A border of a string is a proper prefix that is also a
suffix (shorter than the whole string). Given a string, output the lengths
of all its borders in increasing order.}

\begin{patternbox}
\emph{``All borders of a string''} is a direct application of the
\textbf{prefix function} (KMP failure function). The longest border of the
whole string is $\pi[n-1]$; following the ``failure chain''
$\pi[n-1]\to\pi[\pi[n-1]-1]\to\cdots$ enumerates every border.
\end{patternbox}

\subsection*{Prefix function and the failure chain}

The prefix function $\pi[i]$ is the length of the longest proper border of
the prefix $s[0..i]$. A key fact: the set of all border lengths of the
\emph{whole} string is exactly $\{\pi[n-1],\ \pi[\pi[n-1]-1],\ \ldots\}$
down to $0$. This is because the longest border's own longest border is
the next-longest border of the whole string, and so on. Collecting the
chain and reversing it yields all borders in increasing order.

\begin{ideabox}[Borders Nest Like Russian Dolls]
Every border of $s$ is itself bordered by the next-shorter border. So the
failure chain from $\pi[n-1]$ visits all of them in decreasing order --
each obtained by looking up the prefix function at the previous border's
last index.
\end{ideabox}

\subsection*{Algorithm}

\begin{enumerate}
  \item Compute the prefix function $\pi$ of $s$.
  \item Starting from $b=\pi[n-1]$, repeatedly record $b$ and set
        $b=\pi[b-1]$ until $b=0$.
  \item Reverse the collected lengths and output them (increasing order).
\end{enumerate}

\subsection*{Dry run}

\dryrun{$s=$``abcababcab''.}

\begin{itemize}
  \item The whole string's longest border is ``abcab'' (length $5$):
        $\pi[9]=5$.
  \item $\pi[4]=2$ (border ``ab''); $\pi[1]=0$. Chain: $5,2$.
  \item Increasing order: $2,5$ -- the border lengths.
\end{itemize}

\subsection*{C++ implementation}

\begin{lstlisting}
#include <bits/stdc++.h>
using namespace std;

vector<int> prefixFunction(const string& s) {
    int n = s.size();
    vector<int> pi(n, 0);
    for (int i = 1; i < n; i++) {
        int j = pi[i - 1];
        while (j > 0 && s[i] != s[j]) j = pi[j - 1];
        if (s[i] == s[j]) j++;
        pi[i] = j;
    }
    return pi;
}

int main() {
    string s;
    cin >> s;
    int n = s.size();
    vector<int> pi = prefixFunction(s);

    vector<int> borders;
    for (int b = pi[n - 1]; b > 0; b = pi[b - 1]) borders.push_back(b);
    reverse(borders.begin(), borders.end());

    for (int b : borders) cout << b << " ";
    cout << "\n";
    return 0;
}
\end{lstlisting}

\begin{complexitybox}
\textbf{Time Complexity.} $O(n)$ for the prefix function; the failure
chain has $O(n)$ total length. \textbf{Space Complexity.} $O(n)$.
\end{complexitybox}

\begin{takeawaybox}
Borders are the prefix function's reason for existing. Walking the failure
chain from $\pi[n-1]$ lists them all. This chain is also how KMP recovers
after a mismatch and how periods (next section) are derived.
\end{takeawaybox}
